
% 来源：EPUB part044.xhtml
\section{考虑市场很可能是错的}

其实每个公司最后都会成为那个它本该成为的公司，股价也是。没有长期的莫名其妙的股价，但短期会有。（2023-03-25）

% 来源：EPUB part045.xhtml
\subsection{价格围绕价值上下波动}

价值不是一个概念，是真实的东西。（2010-08-14）

所谓的价值指的就是内在价值。马克思说的“价格围绕价值上下波动”里的价值指的也是内在价值。（2010-03-03）

公司的内在价值就是其未来自由现金流的折现，所谓现在公司的净值应该是被包括在未来现金流里面的。未来现金流折现只是一个概念或一种思考方式。一般来讲，任何人只要试图用未来现金流的计算公式去计算公司的内在价值时，就说明其实他还不太懂自己在干什么。（2011-01-27）

公司价值取决于公司生命周期里的总盈利的折现。好公司生命力强，一般会活得好且长。股市长期来说就是称重器，早晚会体现公司价值的。（2022-04-18）

\BookQuestion{网友：都说尊重常识，请举个常识的例子？}

我应该没这么说过，因为常识的定义不是很清楚，是你的常识、我的常识，还是大家的共识？商业模式好的企业股价最终会上来，这是我的常识，但同意的人可能连10\%都不到，不然不会那么多人在投资上赚不到钱的。（2024-06-26）

% 来源：EPUB part046.xhtml
\subsection{波动是朋友}

“根据股票的波动性来判断风险是很傻的。我们认为只有两种风险：一，血本无归；二，回报不足。有些很好的生意也是波动性很大的，比如喜诗糖果通常一年有两个季度都是亏钱的。反倒是有些烂透了的公司生意业绩很稳定。”（芒格）

想想周围有多少人还在用“股票的波动性”来衡量风险吧。（2012-06-26）

\BookQuestion{网友：记得您是习惯用市价盘买卖股票的，不知道您会不会受到任何短期波动的影响？买了以后马上跌5\%肯定也会觉得惋惜吧？}

你要是真的拥有这家公司，你不会在意的，不然因为股价永远在波动，你每时每刻都会“觉得惋惜”。（2016-01-14）

如果真是大多数投资者都在赌的话，那一定会有个时点可以赚大钱。

我是说一定会有个时点（时间长了会有很多时点）价值投资者能赚大钱。这里时点指的是一个时间区域。（2010-04-20）

\BookQuestion{网友：巴菲特曾说，如果市场总是有效的，他只会成为一个在大街上手拎马口铁罐的流浪汉。}

个人理解：如果市场总是有效的，那我赚啥子钱嘛？其实市场确实在很多时候是很有效的。但市场一定会有胡来的时候。（2012-07-28）

呵呵，波动是朋友。（2010-05-10）

其实自己真的明白的公司跌得重是值得高兴的事情，除非你其实不真懂。（2024-09-18）

\BookQuestion{网友：难过没有闲钱加仓了，跌得重但是没有钱加仓了还要开心吗？}

关你什么事！？（2024-09-18）

“关你什么事”的意思是针对“没有闲钱加仓了”而言的。既然没闲钱了，那股价下跌就跟你无关了。（2024-09-27）

嘴里说着熊市牛市啥的，心里其实想的都是赚市场的钱，其实就是投机的意思。投资赚的是公司赚的钱。这两者有着很大的区别。（2023-12-31）

我不对抗市场，我尽量不理睬市场而已。（2024-06-16）

你可能专注于猜市场很多年了。试试以后忘记市场，专注研究企业。那么多人专注于研究市场几十年都退不了休，你也许可以想想为什么？（2023-05-11）

原来已经计划狠狠买些拼多多和腾讯的，但觉得市场很悲观，应该会低迷一段时间，我可以靠多卖点put（看跌期权）来降低买入的成本，结果……我总是说，任何的考虑市场都很可能是错的，现在就是个具体的例子。（2024-10-01）

% 来源：EPUB part047.xhtml
\subsection{任何股票都只有一个真正的买家}

对投资而言，很重要的一点是你能把上市公司看成非上市公司，然后理性地去想应该怎么做。（2012-11-11）

其实会买公司股票的人只有一个，那就是公司自己，别人都是过客。

想通了这个道理，别的都不是问题了。

我讲的是事实。举个例子，假如苹果经营不变，但所有人都要把股票卖出来，那谁会买？用什么钱买？（2020-10-17）

所有的股票都只有一个真正的买家会最终影响到股价，其他人对股价的影响都是浮云。这个真正的买家就是公司自己，买股票的钱来自盈利。（2013-03-02）

未来净现金流就是这个意思啊。（2012-04-19）

其实只要把公司想象成非上市公司，而且只有一个股东，大概就容易明白了？用非上市公司的办法非常好，不然容易忽略本质的东西。（2013-04-20）

想明白以前最好别碰股票，慢慢想吧。（2013-04-19）

“市场有效”指的是“市场”早晚会有效的，这里的市场其实指的不是市场而是企业本身！比如，假如市场就是不喜欢网易，网易的股价一直在几千万市值，一年可以赚20亿美元，然后股价涨了，你觉得这和“市场”有多大关系？（2016-12-16）

再说一个简单的例子：好几年前，我在某个场合碰到蒂姆·库克（Tim Cook）。他问我你干嘛的，我说，投资。他问，投啥了？我说，苹果。他再问，还有呢？我说，没了，基本上都是苹果。然后他没接话，可能他觉得他见的投资人里没这样的。然后我说，如果苹果一直这个价钱，十来年后我可能就是苹果的唯一股东了哈，因为苹果的盈利会用来把除了我以外的所有股票都买回去的，他听了后微笑了一下，啥也没说。大家都知道，很可惜这个事情没能发生，尤其是有个叫巴菲特的人跑进来把我的如意算盘给搅和了。（2020-10-18）

\BookQuestion{网友：投资一个公司，你会不会想别人为什么会低价卖给你，你又为什么可以高价卖出去？这是不是一种投机的心理？}

呵呵，这个问题很有趣。实际上，我买的时候是不考虑是不是有人从我手里买的。我要假设如果它不是个上市企业，这个价钱我还买不买。你如果明白这点了，价值投资的最基本的概念就有了，反之亦然。（2010-04-03）

\BookQuestion{网友：大道做投资时，特别是思考的时候，有没有一种孤独感？}

从来没有过，因为我从来不在乎别人怎么看投资。只是偶尔会想说：you idiot（你个傻子）！（2024-03-14）

% 来源：EPUB part048.xhtml
\subsection{投自己明白的生意，就不容易恐惧和贪婪}

\BookQuestion{网友：请问在投资中，是心态重要，还是专业知识重要？}

我怎么认为不重要。芒格认为理性最重要。我们认为平常心最重要。好像这俩是一个东西。（2010-05-30）

\BookQuestion{网友：要做到平常心好难，是不是再怎么努力也无法完全做到？有什么方法吗？}

股市掉这么点儿就这样了？平常人是很难有平常心的，所以平常心又叫不平常心。如果你关心的不是事物的本质，没有平常心是正常的。倒过来也一样，如果你关注的只是事物的本质，平常心自然就在那里了。

在错的道路上是没办法有平常心的，比如用着margin（保证金）\footnote{用保证金交易是一种以现金或证券作为抵押品向经纪商融资或融券，进而放大交易杠杆的行为。当交易标的呈现不利波动且交易者无法追加保证金时，会被经纪商强制平仓。}的人们，昨天那种情况下，你懂的。（2015-05-29）

\BookQuestion{网友：有两只我一直研究的股票现在比较便宜，但还没到不用计算器都能看出来很便宜的时候。因为怕踏空，经常有想买入的冲动。又担心现在买了，股价到了很便宜的时候又没有子弹了。你如何处理这种情况？}

我很少考虑踏空的问题，也不太担心没子弹。实在觉得便宜就买，没钱了就打球去，没啥可着急的，不然永远都没法解决你那个问题。最重要的是，你是不是真的懂你想投的那个公司，如果还不是的话，那就再等等，多看看想想。（2011-09-29）

\BookQuestion{网友：您在打球的过程中，会不会因为想到输赢而紧张，导致动作变形？}

平常心在这里指的就是只想每一杆怎么打，别想结果，能做到的赢面就大了。（2015-06-28）

会打高尔夫的人大概都知道，凡是想打远点儿的时候都是最容易犯错的时候。（2013-10-20）

\BookQuestion{网友：请问您如何管理自己的心态？}

我不管理心态，因为投资是不需要管理心态的，但投机需要。（2020-10-28）

需要胆量的是投机。（2010-03-25）

投资不需要勇气，也就是说当你需要勇气时你就危险了。（2010-03-16）

\BookQuestion{网友：如何对待大跌当中的“恐惧”情绪？}

跟心理素质没有关系，恐惧是因为你不明白所谓“价值”为何物。你根本就没明白上市公司与不上市公司的区别，所以下跌你还会恐惧。恐惧的程度与了解程度成反比。只要你投的公司是自己明白的生意，就不容易恐惧，也不容易贪婪。（2011-10-25）

\BookQuestion{网友：相信您在买入很多企业后也会出现一段时间内的股价大跌，请问您最惨时亏损了多少？是如何做心理安慰的？}

你如果买的是好公司，没用margin，那你的公司大掉时你应该是高兴的。我买了苹果后，苹果从最高点掉下来超过40\%的好像就有好几次，

20\%\( - \)30\%的就多了去了（我没认真看过，也不在意）。任何一个公司的真正买家就只有它自己，苹果掉了的时候，哪怕我没钱再多买，我也是高兴的，因为苹果这么赚钱，公司一定会继续买的。如果你不太小心地买了家烂公司，安慰自己的办法大概就是不停地对自己说：以后再也不赌了，再也不赌了，再也不赌了。于是从此可以打败90\%以上的“投资人”了。（2024-03-31）

我们拿着苹果拿着茅台这么多年，老是有人在旁边一直说：你好难啊！我难吗？！难在哪里了？什么事情比这个更容易呢？

拿着好公司比爬山这种东西要简单甚至容易很多，如果你能确定拿着的是“好公司”的话。如果你不知道拿着的是不是好公司，那你坚持什么？爬山这种事情对很多人是需要坚持的。拿着好公司的时候为什么会有要坚持的感觉？是因为你不喜欢吗？（2024-08-15）
